\newpage
\section{State Machines and Languages}

\subsection{Languages and Alphabets}

A \emph{language} \(L\) is a subset of words of an \emph{alphabet} \(\Sigma\) which is a finite 
non-empty set of symbols which always includes the \(\varepsilon\)-word, the empty word.

\subsubsection{Concatenation}

An operation which can be done with members of an alphabet is the so called \emph{concatenation} which is like 
basic addition. Therefore, it is associative and when added \(\varepsilon\), the zero of alphabets it odes not affect the words 
or its length.

\subsection{Inductive Definitions}

In the field of languages and alphabets an inductive function \(f: \Sigma^* \to X\) will be defined 
in the following way.

\begin{enumerate}

    \item \emph{The base case}: \(f(\varepsilon)\)
    
    \item \emph{The induction step}: where the actual behavior of the function is defined. Helper 
          functions can also be defined using the same structure.

\end{enumerate}

It also has the following properties 

\begin{itemize}
    
    \item \(|\varepsilon| = 0\)
    
    \item For \(w = v + a\) then \(|w| = |v| + 1\) for \(a \ne \varepsilon\)

\end{itemize}

\textbf{Example:}

Define a function which count the occurrences of \(ue\) in a word with letters of the alphabet 
\(\Sigma = \{a, \dots, z\}\).

\begin{align*}
    f(\varepsilon) &= 0 \\
    f(w \cdot x) &=
    \begin{cases}
        f(w) & \text{if } x \ne e \\
        \text{help}(w) & \text{if } x = e
    \end{cases}
\end{align*}

The auxiliary function \( \text{help}: \Sigma^* \to \mathbb{N} \) is defined as:

\[
    \text{help}(w \cdot x) =
    \begin{cases}
        f(w) + 1 & \text{if } x = u \\
        f(w)     & \text{otherwise}
    \end{cases}
\]

That is, \( \text{help}(w) = f(w) + 1 \) if the last letter of \( w \) is \( u \), and \( f(w) \) otherwise.

\subsubsection{Head and Tail}

Head returns the first symbol of a word and is defined by 

\begin{align*}
    \operatorname{head}(\varepsilon) = \varepsilon \\ 
    \operatorname{head}(\varepsilon + a) = a \\ 
    \operatorname{head}(v + a) = \operatorname{head}(v) 
\end{align*}

Tail removes the first symbol of a word and is defined by 

\begin{align*}
    \operatorname{tail}(\varepsilon) = \varepsilon \\ 
    \operatorname{tail}(\varepsilon + a) = \varepsilon \\ 
    \operatorname{tail}(v + a) = \operatorname{tail}(v) + a  
\end{align*}


\subsection{Fixes}

Let \(\Sigma\) be an alphabet and let \(v, w \in \Sigma^{*}\) be words.  

\begin{itemize}

    \item \(v\) is called a \emph{substring} of \(w\) iff \(\exists s, t \in \Sigma^{*} : w = s v t\).  

    \item \(v\) is called a \emph{suffix} of \(w\) iff \(\exists s \in \Sigma^{*} : w = s v\) - we write \(v \sqsupseteq w\).  

    \item \(v\) is called a \emph{prefix} of \(w\) iff \(\exists t \in \Sigma^{*} : w = v t\) - we write \(v \sqsubseteq w\).  

    \item \(v\) is called a \emph{proper substring (prefix/suffix)} of \(w\) iff \(v \neq \varepsilon \wedge v \neq w\) and \(v\) is a substring (prefix/suffix) of \(w\).  
    We then write \(v \sqsubset w\) or \(v \sqsupset w\).  

\end{itemize}

\(\varepsilon\) is prefix suffix and substring, of all words. 

\subsection{The Canonical Order}

Let \(\Sigma = \{a_{1}, \dots, a_{n}\}\) be an alphabet with \(n \geq 1\), and let 
\(< \subseteq \Sigma \times \Sigma\) be an order on \(\Sigma\) with
\(a_{1} < a_{2} < \dots < a_{n}\).  

We define the \emph{canonical order} on \(\Sigma^{*}\) as follows:  

\[
    u < v 
    \quad \text{iff} \quad 
    \lvert u \rvert < \lvert v \rvert 
    \;\;\vee\;\; 
    \bigl(\lvert u \rvert = \lvert v \rvert \;\wedge\; 
    u = w \cdot a_{i} \cdot u' \;\wedge\; v = w \cdot a_{j} \cdot v' 
    \;\wedge\; i < j \bigr)
\]

for arbitrary \(w, u, u', v, v' \in \Sigma^{*}\).  

\subsection{Languages}

Let \(\Sigma\) be an alphabet. A \emph{language} \(L\) over \(\Sigma\) is a subset of \(\Sigma^{*}\) (\(L \subseteq \Sigma^{*}\)).  

\[
    \mathcal{P}(\Sigma^{*}) = \{ L \mid L \subseteq \Sigma^{*} \},
\]

the power set of \(\Sigma^{*}\), is the set of all languages over \(\Sigma\).  

Furthermore, \(L_{\emptyset} = \emptyset\) and \(L_{\varepsilon} = \{\varepsilon\}\).  

Another parts of a language is the \emph{syntax} which defines the words in the language, and the  \emph{semantic} which determines 
the meaning of words. 

\subsection{Pumping Lemma}

For a language \(L\) which can be recognized by a state machine, there is a natural number \(n\) such 
that for each word in \(L\) with length greater than or equal to \(n\), the word \(w\) can be decomposed 
into three subwords \(x, y, z\).

\begin{itemize}
    
    \item \(|xy| \leq n\)
    
    \item \(y \neq \varepsilon\)
    
    \item \(xy^i z \in L\) for all \(i \geq 0\)

\end{itemize}

\textbf{Example:}

Let \(L = \{ a^n b^n \mid n \geq 0 \}\). This language is \textbf{not} regular, and the Pumping 
Lemma can be used to prove this by contradiction.

\subsection{Algorithm}

Given two alphabet \(\Sigma_1, \Sigma_2\) and \(\mathcal{D} \subseteq \Sigma_{1}^{*}\). An \emph{algorithm} is a total function 

\[
    A: \mathcal{D} \to \Sigma_{2}^{*}
\]

Two algorithms are equivalent if and only if 

\[
    \forall x \in \mathcal{D} : A(x) = B(x)
\]

\subsection{The Decision Problem}

The \emph{decision problem} \(\langle \Sigma, L \rangle\) for a given alphabet \(\Sigma\) and a language \(L \subseteq \Sigma^{*}\) is to decide, for each \(w \in \Sigma^{*}\), whether  

\[
    w \in L 
    \quad \text{or} \quad 
    w \notin L .
\]

Therefore, one also speaks of the \emph{word problem} \(\langle \Sigma, L \rangle\).  

An algorithm \(A \colon \Sigma^{*} \to \{0, 1\}\) \emph{solves} the decision problem \(\langle \Sigma, L \rangle\) if, for all \(w \in \Sigma^{*}\), we have

\[
    A(w) =
    \begin{cases}
    1 & \text{if } w \in L , \\
    0 & \text{if } w \notin L .
    \end{cases}
\]

We also say: \(A\) \emph{recognizes} \(L\).  

\subsection{Recursive Languages}

Let \(\Sigma\) be an alphabet and let \(L \subseteq \Sigma^{*}\) be a language.  
\(L\) is called \emph{recursive} if and only if there exists an algorithm \(A\) that solves the decision problem \(\langle \Sigma, L \rangle\).  

\subsection{Enumeration Problem}

Let \(\Sigma\) be an alphabet, \(L \subseteq \Sigma^{*}\) a language, and \(n \in \mathbb{N}\).  
The \emph{enumeration problem} \(\langle \Sigma, L, n \rangle\) for \(L\) consists of listing the canonical first \(n\) words of \(L\).  

An algorithm \(A \colon N \to \mathcal{P}(L)\) is an \emph{enumeration algorithm} for \(L\) if, for every input \(n \in N\), \(A\) outputs 
the canonical first \(n\) words \(w_{0}, w_{1}, \dots, w_{n-1}\) of the language \(L\).  

\subsection{Grammars}

Let 

\[
    G = (V, \Sigma, R, S)
\]

where:

\begin{itemize}

    \item \( V \) is a finite set of \emph{variables} (also called non-terminal symbols).
    
    \item \( \Sigma \) is a finite set of \emph{terminal symbols}, such that \( V \cap \Sigma = \emptyset \).
    
    \item \( R \) is a finite set of \emph{production rules} of the form 
    
        \[
            A \rightarrow \alpha,
        \]

        where \( A \in V \) and \( \alpha \in (V \cup \Sigma)^* \).
    
    \item \( S \in V \) is the \textbf{start symbol}.

\end{itemize}

Thus, a grammar \( G \) defines a \emph{formal language} \( L(G) \) as:

\[
    L(G) = \{\, w \in \Sigma^* \mid S \Rightarrow^* w \,\}
\]

where \( \Rightarrow^* \) denotes the reflexive and transitive closure of the derivation relation defined by \( R \).

\subsubsection{Context Free Grammars}



